% Requires pdfpages. Compile from repository root or define MapAtlasDir.
% fitpaper preserves verified A4/A3 page sizes; do not rescale to textwidth.
\providecommand{\MapAtlasDir}{research/2026-a/visuals/cartography}
\includepdf[pages={1},fitpaper=true,pagecommand={}]{\MapAtlasDir/print_ready_maps.pdf}
\includepdf[pages={2},fitpaper=true,pagecommand={}]{\MapAtlasDir/print_ready_maps.pdf}
\includepdf[pages={3},fitpaper=true,pagecommand={}]{\MapAtlasDir/print_ready_maps.pdf}
\includepdf[pages={4},fitpaper=true,pagecommand={}]{\MapAtlasDir/print_ready_maps.pdf}
\includepdf[pages={5},fitpaper=true,pagecommand={}]{\MapAtlasDir/print_ready_maps.pdf}
\includepdf[pages={6},fitpaper=true,pagecommand={}]{\MapAtlasDir/print_ready_maps.pdf}
\includepdf[pages={7},fitpaper=true,pagecommand={}]{\MapAtlasDir/print_ready_maps.pdf}
\clearpage
\section*{Map captions}
\normalsize
\noindent\textbf{Map 1.} Geographical distribution of the 25 candidate cities. Locator points follow the archived urban-area index. Coral markers identify the four accessibility case studies; the Europe inset separates nearby labels. Basemap: Natural Earth; world projection: Robinson.\par\medskip
\noindent\textbf{Map 2.} Learning-facility accessibility in Copenhagen. Within the retained populated grids, 99.4\% of residents have modelled mean walking times to the nearest 20 learning POIs at or below 30 minutes. Colours show five time classes; the inset uses its own scale. Sources: archived Bruno et al. / SONY CSL accessibility data, OpenStreetMap and Natural Earth.\par\medskip
\noindent\textbf{Map 3.} Learning-facility accessibility in Toronto. Within the retained populated grids, 98.2\% of residents have modelled mean walking times to the nearest 20 learning POIs at or below 30 minutes. Colours show five time classes; the inset uses its own scale. Sources: archived Bruno et al. / SONY CSL accessibility data, OpenStreetMap and Natural Earth.\par\medskip
\noindent\textbf{Map 4.} Learning-facility accessibility in Dubai. Within the retained populated grids, 43.7\% of residents have modelled mean walking times to the nearest 20 learning POIs at or below 30 minutes. Colours show five time classes; the inset uses its own scale. Sources: archived Bruno et al. / SONY CSL accessibility data, OpenStreetMap and Natural Earth.\par\medskip
\noindent\textbf{Map 5.} Learning-facility accessibility in Shanghai. Within the retained populated grids, 58.2\% of residents have modelled mean walking times to the nearest 20 learning POIs at or below 30 minutes. Colours show five time classes; the inset uses its own scale. Sources: archived Bruno et al. / SONY CSL accessibility data, OpenStreetMap and Natural Earth.\par\medskip
\noindent\textbf{Map 6.} Four-city comparison of learning-facility accessibility. Identical walking-time classes support comparison of spatial patterns; each city has its own scale bar. The population-weighted shares use the retained urban population grids. Sources: Bruno et al. / SONY CSL, OpenStreetMap and Natural Earth.\par\medskip
\noindent\textbf{Map 7.} Satellite context of the four accessibility case studies. NASA HLS S30 RGB imagery covers the same geographical extents as the accessibility panels. Acquisition dates and city-specific scales appear on the maps. The source product has 30 m pixels; the displayed images are NASA GIBS browse renderings.\par\medskip
